% !TEX program = xelatex
% Шаблон научной статьи НИТУ «МИСиС»
% Заполните блоки ниже и скомпилируйте: xelatex article.tex

\documentclass{misisarticle}

% --- Классификационные коды (оставьте пустым {}, если не требуются) ---
\grnti{28.23.15}
\asjc{1702}
\oecd{1.02, 3.02}

% --- Название ---
\title{Классификация заболеваний сетчатки по снимкам глазного дна с помощью методов глубокого обучения: сравнение свёрточных и трансформерных архитектур в условиях дисбаланса классов }

% --- Авторы: \misisauthor{ФИО}{аффилиация}{город, страна}{email} ---
\misisauthor{А.С.Бабенко}{кафедра инженерной кибернетики НИТУ «МИСиС»}{Москва, Россия}{m2304080@edu.misis.ru}
\misisauthor{A.C.Кожаринов}{кафедра инженерной кибернетики НИТУ «МИСиС»}{Москва, Россия}{kozharinov@misis.ru}

% --- Аннотация и ключевые слова ---
\abstract{%
  В работе рассматривается задача автоматической классификации заболеваний сетчатки по цветным снимкам глазного дна. Качество существующих моделей снижается при дисбалансе классов и при переносе на данные с другим оборудованием и другой популяцией пациентов. Цель работы - повысить точность и устойчивость классификации. Для этого исследуются предобученная архитектура ResNet-18, реализуемая с помощью PyTorch и TorchVision. В работе используется предобработка снимков (нормализация освещённости, кадрирование области диска, восстановление некачественных снимков) и приведение их к единому формату, также применяется аугментация и функции потерь, учитывающие дисбаланс классов (focal loss, взвешенная кросс-энтропия). Модели обучаются на открытых наборах данных. Качество проверяется с помощью метрик ROC-AUC, macro-F1, Accuracy и Balanced Accuracy.
}

\keywords{%
  снимки глазного дна, сверточные нейронные сети, диабетическая ретинопатия, дисбаланс классов, классификация заболеваний сетчатки, обработка изображений
}

\begin{document}

\maketitle

% ============================================================
\section{Введение}
Цветные снимки глазного дна позволяют исследовать изменения сетчатки и используются при разработке систем автоматического распознавания офтальмологических заболеваний. Открытые данные для диабетической ретинопатии, например IDRiD, содержат изображения и специализированную разметку, пригодные для обучения таких систем~\cite{idrid}. Однако результаты, полученные для распространённой патологии, нельзя непосредственно переносить на редкую опухоль: меняются диагностические признаки, состав пациентов и условия получения изображений.

В настоящей работе рассматривается ретинобластома --- злокачественная внутриглазная опухоль детского возраста. Исследования автоматического анализа её изображений уже существуют~\cite{kaliki2023}, поэтому сам факт применения нейронной сети не является новым. Существенным ограничением остаётся доступность разнообразных размеченных данных для внутриглазных опухолей~\cite{sun2025}. При небольшом числе независимых наблюдений модель может запоминать особенности отдельных изображений, а при различии источников --- использовать признаки съёмки вместо признаков заболевания. Аугментация --- создание преобразованных вариантов исходных снимков --- увеличивает число обучающих примеров, но не число независимых пациентов.

Цель исследования --- оценить, улучшает ли промежуточное обучение ResNet-18 на офтальмологических изображениях распознавание ретинобластомы при ограниченной целевой разметке. Здесь перенос обучения означает использование параметров, полученных на другой задаче, как начального состояния модели. Исследовательский вопрос состоит в том, полезен ли дополнительный этап между предобучением на ImageNet и обучением на ретинобластоме. Основная гипотеза: промежуточная классификация степени диабетической ретинопатии на IDRiD повышает macro-F1 относительно прямого переноса ImageNet при одинаковых целевых данных, разбиениях и процедуре настройки. Основное сравнение проводится с обычной кросс-энтропией, чтобы отделить эффект переноса от эффекта функции потерь. Чувствительность к ретинобластоме анализируется отдельно: рост средней метрики не должен скрывать увеличение числа пропущенных положительных примеров.

Планируемый вклад --- сравнение двух схем переноса при фиксированной архитектуре, оценка влияния функций потерь и зависимости качества от объёма целевой разметки. Проверка данных и устранение дубликатов являются условиями достоверности этого сравнения. Разработка новой архитектуры и доказательство клинической эффективности в задачи работы не входят.



% ============================================================
\section{Основная часть}

\section{Анализ работ}
\subsection{Исследования по предметной задаче}
В работе Kaliki и соавторов~\cite{kaliki2023} исследуется обнаружение ретинобластомы и отнесение поражённых глаз к клиническим группам по фотографиям Retcam II. Авторы используют детектор на основе MobileNet V2 SSD, выделяют признаки опухоли и окружающих структур и обучают классификатор; отрицательные изображения относятся к непоражённым глазам пациентов с односторонней опухолью. Это прямой предметный аналог, но его клинические данные и состав отрицательной группы отличаются от произвольного набора изображений из Интернета. Представленное в статье разбиение изображений само по себе не устанавливает независимость пациентов между обучением и тестом.

Sun и соавторы~\cite{sun2025} публикуют набор ультраширокопольных снимков с нормой и пятью видами внутриглазных опухолей, включая ретинобластому. В отличие от бинарной задачи с одними здоровыми контролями, эти данные допускают изучение различения опухолевых состояний. В наборе имеются повторные снимки одних глаз; авторы описывают стратификацию по изображениям, что не гарантирует разделения пациентов. Кроме того, отсутствуют некоторые состояния, похожие на опухоли. Ультраширокопольная съёмка отличается от обычной по полю зрения и способу получения изображения. Поэтому набор рассматривается как кандидат для целевой задачи, а не как доказательство переносимости между условиями съёмки.

Бурсов и Сафонова~\cite{bursov2025} исследуют дублированные и модифицированные изображения в открытых наборах по диабетической ретинопатии. Их работа относится к качеству данных, а не к построению классификатора ретинобластомы. Для предлагаемого эксперимента из неё следует необходимость проверки пересечений как внутри целевого набора, так и между источниками промежуточного обучения.

Комаровских и Подтынных~\cite{komarovskikh2023} рассматривают раннюю диагностику первичной открытоугольной глаукомы по клиническим примерам с использованием нейронных сетей. Эта публикация близка по медицинской области, но не является прямым сравнением свёрточных моделей на цветных фотографиях ретинобластомы. Её диагностические показатели нельзя использовать как ожидаемое качество предлагаемой модели.

\subsection{Методологические аналоги}
Васильева, Ширеторова и Гримм~\cite{vasilieva2025} сопоставляют DeepLabV3+, BiSeNetV2 и YOLOv8-seg для сегментации дорожных трещин на объединённых открытых данных, включая Crack500 и CrackForest. Это аналог сравнительного эксперимента на неоднородных изображениях, но его пиксельная задача и объекты отличаются от распознавания опухоли по снимку.

Ерещенко~\cite{ereshchenko2023} рассматривает сиамскую сеть и распознавание дефектов стальных поверхностей при ограниченном числе примеров; обсуждаются NEU и DAGM. Работа полезна как методологический аналог обучения при дефиците данных. Однако перенос выводов между промышленными текстурами и офтальмологическими изображениями требует отдельного эксперимента.

Курочкин, Прун и Балаев~\cite{kurochkin2024} исследуют федеративное обучение классификатора изображений на грид-системе персональных компьютеров с ограничениями обмена частью обучающих данных. Эта работа относится к организации распределённого обучения. Она не подтверждает пользу предметного переноса и не обосновывает включение федеративной схемы в настоящий эксперимент.

Архитектурной основой служит остаточное обучение He и соавторов~\cite{resnet}. Focal loss предложена Lin и соавторами~\cite{focal} для уменьшения влияния легко распознаваемых примеров при обучении детектора. Её применение к классификации ретинобластомы является проверяемым методическим выбором, а не установленным способом улучшения качества.

% Редакторский статус источников вынесен из основного изложения.
% Требуется согласовать, какая публикация удовлетворяет требованию о работе преподавателя.
Работа Федотова, Черноморца и Михелева~\cite{fedotov2018} относится к сегментации сосудов глазного дна. Её библиографические реквизиты и DOI подтверждены регистрационной записью, однако полный текст недоступен для проверки. Поэтому эта позиция сохранена как тематическая ссылка и не используется для обоснования выбора модели или ожидаемого качества.

Таким образом, рассмотренные исследования не отвечают на узкий вопрос настоящей работы: даёт ли дополнительное обучение одной и той же ResNet-18 на офтальмологических данных преимущество перед прямым переносом ImageNet при контроле независимости целевой выборки? Для ответа необходимы раздельное варьирование маршрута переноса и функции потерь, единые разбиения и оценка редкого класса.

\section{Постановка задачи}
\subsection{Данные и границы постановки}
В исходном плане указан Retinoblastoma Dataset пользователя kedimestan на Kaggle~\cite{rbkaggle}. Его описание ссылается на проект Roboflow~\cite{rbroboflow}, где перечислены метки \texttt{healthy}, \texttt{RB}, \texttt{RL}, \texttt{RR} и \texttt{Unlabeled}. Развёрнутое описание проекта отсутствует. Наличие метки \texttt{healthy} подтверждено для первоисточника, но не устанавливает наличие соответствующих файлов в конкретной копии Kaggle. Значения \texttt{RL} и \texttt{RR} не расшифрованы; относить их к отрицательному или положительному классу по одному названию нельзя. \texttt{Unlabeled} не является диагностической меткой.

При проверке доступного архива Kaggle версии~1 обнаружены JPEG-файлы в общей директории, без отдельного файла меток и классовых подкаталогов. Имена содержат обозначения \texttt{RB} и \texttt{aug}; просмотренная выборка включает внешние фотографии глаз и лица, а не изображения глазного дна. Визуально просмотрены девять файлов без префикса \texttt{aug}, равномерно выбранных из списка таких файлов. Это ограниченная проверка содержимого, а не полный аудит всех изображений или версий первоисточника. Число независимых пациентов, клиническое подтверждение меток и отрицательные примеры по этому архиву не установлены. Поэтому он не принимается как готовая выборка для заявленной задачи.

Дальнейшая постановка является \emph{планируемой} и требует целевого набора снимков глазного дна с проверяемыми положительными и отрицательными метками. Тип визуализации, включая обычную цветную фотографию или псевдоцветное изображение сканирующего лазерного офтальмоскопа, должен быть указан для выбранного источника. Возможная замена --- набор Sun и соавторов~\cite{sun2025}. До формирования выборки необходимо установить, означает ли RB наличие видимой опухоли на конкретном снимке или диагноз пациента, в том числе после лечения. Эти типы меток не считаются взаимозаменяемыми. Для первичного бинарного эксперимента предлагается сравнение RB с Normal; включение остальных опухолей в отрицательный класс оформляется отдельным расширением задачи после проверки разметки. Результат RB против Normal не интерпретируется как способность отличать ретинобластому от всех заболеваний. Вспомогательные наборы по глаукоме и ретинопатии не становятся отрицательными примерами автоматически.

\subsection{Вход и выход модели}
Пусть после аудита сформирована выборка
\begin{equation}
 \mathcal D=\{(x_i,y_i,g_i,s_i)\}_{i=1}^{N},
 \label{eq:data}
\end{equation}
где $x_i$ --- трёхканальное изображение глазного дна в формате RGB, $y_i\in\{0,1\}$ --- проверенная метка (1 --- ретинобластома, 0 --- согласованный отрицательный класс), $g_i$ --- идентификатор пациента, $s_i$ --- источник изображения. $N$ обозначает число отобранных снимков до онлайн-аугментации и не задаётся заранее. При отсутствии $g_i$ это явно фиксируется как ограничение, а не заменяется номером файла.

Модель выдаёт оценку $p_\theta(x)\in[0,1]$ принадлежности положительному классу и решение
\begin{equation}
 \widehat y=\mathbf{1}\{p_\theta(x)\geq\tau\},
 \label{eq:decision}
\end{equation}
где $\theta$ --- обучаемые параметры, $\tau$ --- порог, выбираемый только на валидации. Выход сети без отдельной проверки калибровки не трактуется как клиническая вероятность заболевания. Единица предсказания --- снимок. Группировка по пациентам предотвращает попадание снимков одного человека в разные части выборки; сама по себе она не устраняет смещения, связанные с источником данных. Метрики по всем снимкам дают больший вклад пациентам с большим числом изображений, поэтому распределение числа снимков на пациента также включается в отчёт.

\subsection{Критерии качества}
Основной критерий сравнения --- macro-F1:
\begin{equation}
 F_{1,c}=\frac{2TP_c}{2TP_c+FP_c+FN_c},\qquad
 F_{1,\mathrm{macro}}=\frac{F_{1,0}+F_{1,1}}{2}.
 \label{eq:macrof1}
\end{equation}
Здесь $TP_c$, $FP_c$, $FN_c$ --- числа истинно положительных, ложно положительных и ложно отрицательных предсказаний для класса $c$ в схеме «класс против остальных». Равные веса классов уменьшают доминирование многочисленного класса, характерное для доли правильных ответов (Accuracy), но не делают macro-F1 независимой от состава выборки. F1 учитывает одновременно полноту и точность положительных предсказаний. Если модель не предсказывает класс, его F1 принимается равным нулю; разбиение без истинных представителей одного из классов не используется для итогового сравнения.

Для класса ретинобластомы отдельно вычисляются
\begin{equation}
 P_1=\frac{TP}{TP+FP},\quad
 R_1=\frac{TP}{TP+FN},\quad
 S_0=\frac{TN}{TN+FP},
 \label{eq:pr}
\end{equation}
где $TP$, $FP$, $FN$, $TN$ определены относительно $y=1$; $P_1$ --- точность положительных предсказаний, $R_1$ --- чувствительность (полнота), $S_0$ --- специфичность. При отсутствии положительных предсказаний $P_1$ помечается как неопределённая, а $F_{1,1}$ равна нулю при наличии истинных положительных объектов. Дополнительно приводятся $F_{1,1}$, матрица ошибок и средняя точность (average precision, AP):
\begin{equation}
 \mathrm{AP}=\sum_{j=1}^{J}(R^{(j)}-R^{(j-1)})P^{(j)},\quad R^{(0)}=0.
 \label{eq:ap}
\end{equation}
Здесь $P^{(j)}$ и $R^{(j)}$ --- точность положительных предсказаний и полнота после добавления очередной группы снимков с одинаковой оценкой модели при движении от высоких оценок к низким, $J$ --- число различных оценок. Такая запись задаёт AP, а не трапецеидальное интегрирование кривой precision--recall. AP характеризует ранжирование без выбора одного порога; вместе с ней сообщается доля положительных наблюдений, от которой зависит её интерпретация~\cite{apdocs}. Площадь под ROC-кривой (ROC-AUC), сбалансированная точность $\mathrm{BA}=(R_1+S_0)/2$ и Accuracy служат дополнительными характеристиками. Macro-F1 не учитывает повышенную цену пропуска опухоли, поэтому отдельно сообщаются $R_1$ и число $FN$, даже если средний показатель вырос.

Редкость заболевания в популяции не означает, что положительный класс малочислен в отобранном наборе. Фактическое соотношение классов определяется после аудита. Выбор взвешивания обосновывается этим соотношением, а не только медицинской редкостью. Для оценки дефицита разметки планируются кривые качества при уменьшении числа обучающих пациентов с неизменной оценочной выборкой.

\section{Описание методов}
Далее описан предлагаемый протокол. Обучение моделей по нему ещё не выполнено; указанные процедуры являются планом, а не отчётом о результатах.
\subsection{Базовая модель и предметный перенос}
Используется ResNet-18 из TorchVision с фиксированными весами \texttt{IMAGENET1K\_V1}~\cite{torchvision}. Остаточные блоки добавляют преобразование признаков к входному представлению~\cite{resnet}. После глобального усреднения модель формирует вектор $h_\phi(x)\in\mathbb R^{512}$. Исходный классификатор заменяется линейным слоем с двумя выходами:
\begin{equation}
 z=Wh_\phi(x)+b,\qquad
 p_c(x)=\frac{\exp z_c}{\sum_{k=0}^{1}\exp z_k},
 \label{eq:head}
\end{equation}
где $W\in\mathbb R^{2\times512}$, $b\in\mathbb R^2$, $\phi$ --- параметры энкодера (сети, извлекающей признаки изображения), $\theta=(\phi,W,b)$; в формуле~\eqref{eq:decision} используется $p_1(x)$.

Базовый маршрут A: ImageNet $\rightarrow$ целевая задача. Сначала обучается новый классификатор с замороженным энкодером, затем выполняется дообучение всей сети с меньшим шагом для энкодера. Маршрут B: ImageNet $\rightarrow$ вспомогательная офтальмологическая задача $\rightarrow$ целевая задача. После промежуточного этапа вспомогательный классификатор удаляется, создаётся такой же новый бинарный слой и применяется та же процедура целевого обучения. Это позволяет оценить влияние промежуточного представления при неизменной архитектуре.

Для основного варианта B предлагается IDRiD и классификация степени диабетической ретинопатии. На промежуточном этапе используется обычная многоклассовая кросс-энтропия; модель выбирается по macro-F1 вспомогательной валидационной части. Один и тот же промежуточный энкодер в рамках запуска служит начальным состоянием для всех целевых функций потерь. Отдельные дополнительные варианты заменяют IDRiD на Eye Diseases Classification либо EyePACS-AIROGS-light-V2. Наборы не объединяются в одну плоскую систему классов: их метки описывают разные задачи. Последовательное обучение на всех трёх наборах не включается в основное сравнение, поскольку смешивало бы влияние источника и порядка обучения.

IDRiD содержит разметку пяти степеней диабетической ретинопатии (0--4) и риска макулярного отёка~\cite{idrid}; в выбранном промежуточном этапе используется только степень ретинопатии, а сегментационные маски и разметка отёка не подменяют целевые метки. Eye Diseases Classification содержит Normal, Diabetic Retinopathy, Cataract и Glaucoma~\cite{eyediseases}; он нужен для проверки переноса от более разнородной четырёхклассовой задачи. В его описании IDRiD указан среди источников, поэтому независимость этих двух наборов предполагать нельзя. Сходство названий классов также не гарантирует одинаковые критерии разметки; перенос касается параметров энкодера, а не диагностических меток. EyePACS-AIROGS-light-V2 содержит классы referable glaucoma (RG) и non-referable glaucoma (NRG), а также готовые train/validation/test~\cite{airogs}. NRG означает отсутствие глаукомы, требующей направления, а не доказанное отсутствие любой патологии или ретинобластомы. Этот набор служит альтернативным источником промежуточного обучения, но не внешним тестом целевой задачи.

\subsection{Предобработка и аугментации}
До обучения проверяются модальность, читаемость, происхождение и согласованность меток. Удаляется только чёрная рамка вокруг поля глазного дна; кадрирование по диску зрительного нерва не применяется, поскольку может исключить опухоль вне этой области. В качестве исходного варианта предлагается масштабирование с сохранением пропорций и дополнением до квадрата $224\times224$ пикселя. Разрешение необходимо проверить на обучающих снимках: уменьшение может скрыть малые очаги. Это проектное решение, а не установленный оптимум. Для всех вариантов используются одинаковые преобразования. Значения RGB переводятся в $[0,1]$ и нормализуются средними $(0.485,0.456,0.406)$ и стандартными отклонениями $(0.229,0.224,0.225)$, соответствующими исходным весам~\cite{torchvision}.

В обучающей части планируются умеренные повороты, горизонтальные отражения и изменения яркости и контраста. Их диапазоны и вероятности в исходных материалах не заданы: они фиксируются после проверки допустимости преобразований на обучающих снимках, до итоговой оценки. При повороте поле зрения не должно обрезаться; цветовые изменения не должны уничтожать диагностические признаки. Случайное вырезание фрагментов и генеративное восстановление снимков не используются. Нечитаемые изображения исключаются по заранее установленным критериям с отчётом об исключениях. Дополнительная коррекция освещения не включается в основное сравнение. При её исследовании выполняется отдельное сравнение с обработкой и без неё при остальных одинаковых условиях. Для валидации и теста преобразования детерминированы. Все новые аугментации создаются после разбиения и только внутри обучения.

\subsection{Функции потерь и настройка}
Пусть $p_{i,y_i}$ --- вероятность истинного класса для $i$-го обучающего объекта, а $n$ --- число объектов обучающей части. Сопоставляются
\begin{align}
 L_{\mathrm{CE}}&=-\frac1n\sum_{i=1}^{n}\log p_{i,y_i},\label{eq:ce}\\
 L_{\mathrm{WCE}}&=-\frac1n\sum_{i=1}^{n}w_{y_i}\log p_{i,y_i},\label{eq:wce}\\
 L_{\mathrm{FL}}&=-\frac1n\sum_{i=1}^{n}w_{y_i}(1-p_{i,y_i})^\gamma
 \log p_{i,y_i}.\label{eq:fl}
\end{align}
Для взвешенных вариантов $w_c=n/(2n_c)$, где $n_c$ --- число исходных обучающих изображений класса $c$ после очистки; веса пересчитываются внутри каждого разбиения. В focal loss множитель $(1-p_{i,y_i})^\gamma$ уменьшает вклад простых примеров~\cite{focal}. Параметр $\gamma>0$ определяет степень этого подавления; диапазон его подбора должен быть установлен до обучения и проверяться только на валидации. При $\gamma=0$ формула~\eqref{eq:fl} совпадает с~\eqref{eq:wce}; это позволяет проверить реализацию предельного случая. В отличие от исходной параметризации focal loss, здесь явно используются те же обратные частоты классов, что и в WCE; выбранная нормировка является частью протокола. Усреднение во всех формулах выполняется по числу примеров; при реализации библиотечные функции должны воспроизводить эту нормировку. Функции вычисляются через устойчивые операции над логитами $z$, чтобы избежать вычисления логарифма численного нуля.

Основное сравнение включает шесть условий (табл.~\ref{tab:design}). Влияние переноса проверяется прежде всего для B/CE против A/CE; влияние потерь --- внутри каждого маршрута. Для двух дополнительных источников предлагаются отдельные варианты с CE, чтобы не расширять без необходимости основную матрицу.
\begin{table}[htbp]
\centering
\caption{План основных вариантов целевого обучения}
\label{tab:design}
\begin{tabular}{lll}
\toprule
Вариант & Промежуточный этап & Потери \\
\midrule
A/CE & отсутствует & CE \\
A/WCE & отсутствует & WCE \\
A/FL & отсутствует & FL \\
B/CE & IDRiD & CE \\
B/WCE & IDRiD & WCE \\
B/FL & IDRiD & FL \\
\bottomrule
\end{tabular}
\end{table}

Повторная выборка объектов миноритарного класса (oversampling) в основной серии не применяется, чтобы не смешивать её влияние с весами потерь. Оптимизатор, шаги обучения, регуляризация, размер мини-пакета, число эпох, правило ранней остановки и бюджет поиска пока не утверждены. Их требуется зафиксировать в протоколе; общие параметры подбираются в одинаковых диапазонах с одинаковым бюджетом целевой настройки. Для FL отдельно учитывается подбор $\gamma$. Ранняя остановка и выбор контрольной точки выполняются по валидационной macro-F1 при пороге $0.5$, затем на той же валидационной части выбирается рабочий порог. Это настройка по валидации, поэтому её метрика не служит итоговой оценкой качества.

Сравнение оценивает пользу добавленного этапа в целом. Оно не отделяет эффект офтальмологического содержания от эффекта дополнительных шагов оптимизации. Для такого вывода нужен отдельный контроль с сопоставимым вычислительным бюджетом; в основной план он не включён. Время промежуточного и целевого обучения учитывается раздельно.

Порог $\tau$ выбирается по максимуму macro-F1 на валидации; кандидаты определяются по различным валидационным оценкам модели. При равенстве выбирается порог, ближайший к $0.5$, а при повторном равенстве --- меньший. Вместе с этим сообщаются метрики при фиксированном $\tau=0.5$. Отдельная оценка чувствительности при заданной специфичности возможна после согласования её требуемого уровня; тест для выбора порога не используется.

\subsection{Разбиение и проверка гипотезы}
Аудит дублей проводится до разбиения по всем задействованным источникам: сравниваются хеши файлов и декодированных пикселей, ищутся визуально близкие изображения, проверяются имена и происхождение готовых аугментаций. Связи по пациенту и происхождению снимка объединяются: все объекты, связанные хотя бы цепочкой таких связей, образуют неделимую группу. Визуальное сходство проверяется вручную, поскольку похожие снимки разных пациентов не обязательно являются дубликатами. Точные копии удаляются; готовые аугментации не включаются в оценочную часть. Если восстановить оригинал невозможно, такая группа исключается из валидации и теста. Все снимки обоих глаз и повторных визитов одного пациента также остаются в одной части.

При достаточном числе независимых пациентов предусматривается разделение на обучающую, валидационную и тестовую части с сохранением групп и по возможности долей классов. Доли разбиения определяются после аудита числа пациентов и фиксируются до обучения; назначать их сейчас без этих сведений необоснованно. Если редких групп недостаточно для устойчивого отложенного теста, до начала моделирования выбирается вложенная групповая кросс-валидация с числом блоков, допускающим присутствие обоих классов в оценочных частях. Внешний блок не участвует в подборе параметров, а порог определяется во внутреннем цикле. При неизвестных пациентах группировка по дублям снижает утечку, но не доказывает независимость наблюдений; результаты помечаются как предварительные.

Целевые валидационные и тестовые снимки, их дубли и доступные снимки тех же пациентов исключаются из промежуточного обучения. Готовые разбиения вспомогательных наборов используются лишь после проверки групп и пересечений. Вспомогательный test не используется для подбора промежуточной модели. Для всех условий используются одинаковые целевые разбиения и согласованный список случайных начальных состояний. Число повторов устанавливается до обучения с учётом ресурсов. В парных запусках совпадают инициализация нового классификатора и порядок целевых обучающих примеров; различается схема получения энкодера. Во время обучения одного классификационного слоя фиксируются также накопленные статистики пакетной нормализации (BatchNorm), иначе энкодер фактически изменяется.

Для каждого запуска сохраняются предсказания по всем тестовым снимкам. Отдельно сообщаются средние метрики и разброс по запускам. На независимом отложенном тесте неопределённость разности оценивается парным кластерным бутстрепом: пациенты выбираются с возвращением, все их снимки включаются вместе, а для обеих моделей используются одни и те же выбранные группы. В каждой повторной выборке вычисляется средняя по согласованным запускам разность метрик; усреднять вероятности разных запусков вместо этого нельзя, поскольку это оценивало бы ансамбль. Повторные выборки без одного из классов не используются; при частом возникновении таких случаев интервал не считается надёжным. Уровень интервала и число повторных выборок фиксируются заранее. Такой интервал условен на выбранных обучающих данных и не отражает всю неопределённость переобучения модели на новой выборке.

При вложенной кросс-валидации отдельно сообщаются результаты внешних блоков и их разброс. Блоки не считаются независимыми экспериментами: их обучающие части пересекаются. Описанный выше бутстреп отложенного теста не переносится автоматически на объединённые предсказания кросс-валидации.

Для основного сравнения вычисляется
\begin{equation}
 \Delta=F_{1,\mathrm{macro}}^{B,\mathrm{CE}}
       -F_{1,\mathrm{macro}}^{A,\mathrm{CE}}.
 \label{eq:delta}
\end{equation}
В формуле~\eqref{eq:delta} метрики усредняются по согласованным запускам; дополнительно приводятся парные разности каждого запуска. Гипотеза получает поддержку на отложенном тесте, если интервал для $\Delta$ лежит выше нуля. Пересечение интервала с нулём означает недостаточность свидетельств преимущества, а не доказанную эквивалентность моделей. Вывод относится к выбранной выборке и протоколу. Остальные сравнения рассматриваются как поисковые; лучший результат не выбирается по тесту.

Для исследования дефицита разметки планируется обучение на вложенных подвыборках пациентов с сохранением обоих классов. Уровни объёма и один основной уровень для проверки гипотезы определяются после аудита и фиксируются до обучения; остальные точки кривой используются для дополнительного анализа. На каждом уровне обе модели получают одни и те же подвыборки, а состав валидации и теста остаётся неизменным. Публикуются число пациентов и снимков каждого класса. Если подбор параметров использует разметку валидации, она учитывается в общем бюджете разметки: кривая характеризует уменьшение именно обучающей части при фиксированном валидационном бюджете.


% ============================================================
\section{Заключение}

Кратко сформулируйте основные выводы работы.

% ============================================================
% Список литературы — вариант 1: вручную
\begin{thebibliography}{99}
\bibitem{idrid} Porwal P., Pachade S., Kamble R., Kokare M., Deshmukh G., Sahasrabuddhe V., Meriaudeau F. Indian Diabetic Retinopathy Image Dataset (IDRiD): A Database for Diabetic Retinopathy Screening Research. \textit{Data}. 2018. Vol.~3, no.~3. Art.~25. DOI: \url{https://doi.org/10.3390/data3030025}.
\bibitem{kaliki2023} Kaliki S., Vempuluru V.~S., Ghose N. et al. Artificial intelligence and machine learning in ocular oncology: Retinoblastoma. \textit{Indian Journal of Ophthalmology}. 2023. Vol.~71, no.~2. P.~424--430. DOI: \url{https://doi.org/10.4103/ijo.IJO_1393_22}.
\bibitem{sun2025} Sun J., Zhao X., Chen S. et al. An ultra-wide-field fundus image dataset for intelligent diagnosis of intraocular tumors. \textit{Scientific Data}. 2025. Vol.~12. Art.~1521. DOI: \url{https://doi.org/10.1038/s41597-025-05864-2}.
\bibitem{bursov2025} Бурсов А.~И., Сафонова Д.~М. Исследование изображений в открытых датасетах глазного дна при диабетической ретинопатии, предназначенных для обучения нейросетевых алгоритмов. \textit{Сибирский журнал клинической и экспериментальной медицины}. 2025. Т.~40, №~1. С.~218--225. DOI: \url{https://doi.org/10.29001/2073-8552-2025-40-1-218-225}.
\bibitem{komarovskikh2023} Комаровских Е.~Н., Подтынных Е.~В. Результаты применения искусственных нейронных сетей для ранней диагностики глаукомы. \textit{Российский офтальмологический журнал}. 2023. Т.~16, №~2. С.~28--32. DOI: \url{https://doi.org/10.21516/2072-0076-2023-16-2-28-32}.
\bibitem{vasilieva2025} Васильева А.~О., Ширеторова И.~А., Гримм М. Сравнение моделей сегментации дорожных трещин на основе современных нейросетевых архитектур. \textit{Искусственный интеллект в промышленных, коммерческих, медицинских и финансовых приложениях}. Вып.~4. М.: НИТУ «МИСИС», 2025. С.~43--49. \url{https://misis.ru/files/-/97ac0ca5f5966e7f7da4a99f1b3bd35b/kik_2025_1.pdf}.
\bibitem{ereshchenko2023} Ерещенко А.~Г. Исследование распознавания производственных дефектов на стальных поверхностях при помощи компьютерного зрения. \textit{Искусственный интеллект в промышленных, коммерческих, медицинских и финансовых приложениях}. Вып.~1. М.: НИТУ «МИСИС», 2023. С.~138--144. \url{https://misis.ru/files/-/1a3a57deb59c39284cba56d0b54bb308/kik_2023.pdf}.
\bibitem{kurochkin2024} Курочкин И.~И., Прун А.~И., Балаев А.~А. Федеративное глубокое обучение для решения задачи классификации изображений на грид-системе из персональных компьютеров. \textit{Физика элементарных частиц и атомного ядра}. 2024. Т.~55, №~3. С.~585. \url{https://pepan.jinr.ru/index.php/Pepan/article/view/903}.
\bibitem{resnet} He K., Zhang X., Ren S., Sun J. Deep Residual Learning for Image Recognition. \textit{Proceedings of CVPR}. 2016. P.~770--778. \url{https://arxiv.org/abs/1512.03385}.
\bibitem{focal} Lin T.-Y., Goyal P., Girshick R., He K., Doll\'{a}r P. Focal Loss for Dense Object Detection. \textit{Proceedings of ICCV}. 2017. P.~2980--2988. \url{https://arxiv.org/abs/1708.02002}.
\bibitem{fedotov2018} Федотов Е.~А., Черноморец Д.~А., Михелев В.~М. Об эффективности метода последовательной сегментации толстых и тонких кровеносных сосудов глазного дна. \textit{Научные ведомости Белгородского государственного университета. Серия: Экономика. Информатика}. 2018. Т.~45, №~2. С.~312--321. DOI: \url{https://doi.org/10.18413/2411-3808-2018-45-2-312-321}. \textit{Библиографические реквизиты сверены с Crossref; полный текст и экспериментальные сведения не проверены.}
\bibitem{rbkaggle} Albayrak M.~K. Retinoblastoma Dataset. Kaggle. Версия~1, 2025. \url{https://www.kaggle.com/datasets/kedimestan/retinoblastoma-dataset}. Описание и архив проверены 07.11.2026.
\bibitem{rbroboflow} new-workspace-jvb3p. Retinoblastoma Dataset. Roboflow Universe. 2022. \url{https://universe.roboflow.com/new-workspace-jvb3p/retinoblastoma-4fu8f}. Дата обращения: 07.11.2026.
\bibitem{torchvision} TorchVision. ResNet-18: model and pretrained weights documentation. \url{https://docs.pytorch.org/vision/stable/models/generated/torchvision.models.resnet18.html}. Дата обращения: 07.11.2026.
\bibitem{eyediseases} Gunavenkat Doddi. Eye Diseases Classification. Kaggle. \url{https://www.kaggle.com/datasets/gunavenkatdoddi/eye-diseases-classification}. Описание проверено 07.11.2026.
\bibitem{airogs} Kiefer R. EyePACS-AIROGS-light-V2. Kaggle, 2024. DOI: \url{https://doi.org/10.34740/KAGGLE/DSV/7802508}. \url{https://www.kaggle.com/datasets/deathtrooper/glaucoma-dataset-eyepacs-airogs-light-v2}. Описание проверено 07.11.2026.
\bibitem{apdocs} Scikit-learn. \texttt{average\_precision\_score}: documentation. \url{https://scikit-learn.org/stable/modules/generated/sklearn.metrics.average_precision_score.html}. Дата обращения: 07.11.2026.
\end{thebibliography}
\end{document}